%%%%%%%%%%%%%%%%%%%%%%%%%%%%%%%%%%%%%%%%%%%%%%%%%%%%%%%%%%%%%%%%%%%%%%%%%%%%%%%
%%  Chapter 16 — Cross-Domain Unification I:
%%               Cognitive Science, Psychology, Neuroscience,
%%               Interface Theory, and Biology
%%
%%  DEPENDENCIES: Chapters 3–15 (Full CIIR/CRS/Emergent Physics/Audit)
%%
%%  PURPOSE: Derive cognitive, psychological, neuroscientific, and biological
%%           phenomena as projections, constraints, and substructures of CIIR.
%%           Each domain emerges from the primitive triple (S, C, D) and the
%%           operator algebra B(H_C) without introducing new primitives.
%%%%%%%%%%%%%%%%%%%%%%%%%%%%%%%%%%%%%%%%%%%%%%%%%%%%%%%%%%%%%%%%%%%%%%%%%%%%%%%

\chapter{Cross-Domain Unification~I: Mind, Brain, and Life}
\label{ch:cross-domain-I}

\vspace{0.5em}
\noindent
This chapter demonstrates that cognitive science, psychology, neuroscience,
interface theory, and systems biology each emerge as a constrained projection
of the CIIR framework.  Every construction below uses only:
\begin{enumerate}
  \item The reality space $\CR$ and constraint manifold $(\CM, g, \kappa)$
        (Definitions~\ref{def:3.1}, \ref{def:3.7}).
  \item The cognitive Hilbert space $\HC$ and interface map $\Phi$
        (Definitions~\ref{def:3.11}, \ref{def:3.18}).
  \item The operator algebra $\CB(\HC)$ and constraint operators $\hatC_i$
        (Chapter~5).
  \item The CIIR master equation and Lindblad dynamics (Chapter~7).
  \item The primitive triple $(\CS, \CC, \CD)$ (Chapter~14).
\end{enumerate}

No new axioms or primitives are introduced.

% ============================================================
\section{Generic Domain-Embedding Apparatus}
\label{sec:ch16:apparatus}
% ============================================================

Before treating specific domains, we establish the universal embedding
mechanism by which any domain $D$ is realised within CIIR.

\begin{definition}[Domain Projection]\label{def:16.1}
A \emph{domain projection} is a triple $(P_D, \mathcal{O}_D, \Omega_D)$
where:
\begin{enumerate}
  \item $P_D : \HC \to \HC_D$ is an orthogonal projection onto a
        closed subspace $\HC_D \subseteq \HC$ (the \emph{domain
        subspace}).
  \item $\mathcal{O}_D := P_D \CB(\HC) P_D \cong \CB(\HC_D)$ is the
        \emph{restricted operator algebra} (a $C^*$-subalgebra).
  \item $\Omega_D := \{ \rho_D = P_D \rho P_D / \Tr(P_D \rho P_D)
        : \rho \in \mathfrak{D}(\HC),\; \Tr(P_D \rho P_D) > 0 \}$
        is the \emph{domain state space}.
\end{enumerate}
\end{definition}

\begin{lemma}[Domain CPTP Reduction]\label{lem:16.1}
Let $\CL$ be the CIIR Lindblad generator (Definition~7.1).  Define
$\CL_D(\cdot) := P_D \CL(P_D \,\cdot\, P_D) P_D$.  Then $\CL_D$
generates a completely positive, trace-non-increasing semigroup on
$\mathcal{T}_1(\HC_D)$.  If the subspace $\HC_D$ is $\CL$-invariant,
$\CL_D$ is trace-preserving and generates a CPTP semigroup.
\end{lemma}

\begin{proof}
$P_D \CL(P_D \rho_D P_D) P_D$ is a composition of CP maps (sandwich
by $P_D$) and the CP generator $\CL$.  Composition of CP maps is CP.
Trace non-increase follows from $\Tr(P_D A P_D) \leq \Tr(A)$ for
$A \geq 0$.  If $\CL(P_D \HC) \subseteq P_D \HC$, then
$\Tr(\CL_D(\rho_D)) = \Tr(\CL(P_D \rho_D P_D)) = 0$, giving
trace preservation.
\end{proof}

\begin{definition}[Domain Formal Embedding]\label{def:16.2}
A scientific domain $D$ is \emph{formally embedded in CIIR} if there
exists a domain projection $(P_D, \mathcal{O}_D, \Omega_D)$ and a
map
\[
  f_D : (\CM, \CB(\HC), \Phi) \longmapsto (P_D, \mathcal{O}_D, \Omega_D)
\]
such that:
\begin{enumerate}
  \item[(E1)] The states of $D$ are elements of $\Omega_D$.
  \item[(E2)] The operators of $D$ are elements of $\mathcal{O}_D$.
  \item[(E3)] The observables of $D$ are self-adjoint elements of
              $\mathcal{O}_D$.
  \item[(E4)] The dynamics of $D$ are generated by $\CL_D$.
\end{enumerate}
\end{definition}

% ============================================================
\section{Domain I: Cognitive Science and Psychology}
\label{sec:ch16:cogsci}
% ============================================================

\subsection{Formal Definition}

\begin{definition}[Cognitive Domain]\label{def:16.3}
The \emph{cognitive domain} is:
\[
  D_{\mathrm{cog}} := (P_{\mathrm{cog}},\; \mathcal{O}_{\mathrm{cog}},\;
                       \Omega_{\mathrm{cog}})
\]
where $P_{\mathrm{cog}}$ projects onto the subspace
$\HC_{\mathrm{cog}} \subseteq \HC$ spanned by the eigenstates of
the constraint operator $\hatC$ with eigenvalues in a bounded
interval $[\lambda_{\min}, \lambda_{\max}]$:
\[
  \HC_{\mathrm{cog}} := \bigoplus_{\lambda_j \in [\lambda_{\min},\lambda_{\max}]}
                         \ker(\hatC - \lambda_j I).
\]
\end{definition}

\begin{remark}
The interval $[\lambda_{\min}, \lambda_{\max}]$ corresponds to the
\emph{cognitive bandwidth} — the range of constraint intensities that
a given cognitive system can process.  This is not a new primitive; it
is determined by the spectral structure of $\hatC$ (Theorem~5.2).
\end{remark}

\subsection{Ontological Mapping}

\begin{center}
\begin{tabular}{lll}
\toprule
\textbf{Cognitive Concept} & \textbf{CIIR Object} & \textbf{Reference} \\
\midrule
Mental state & $\rho_{\mathrm{cog}} \in \Omega_{\mathrm{cog}}$ & Def.~\ref{def:16.1} \\
Cognitive process & $O \in \mathcal{O}_{\mathrm{cog}}$ & Def.~\ref{def:16.1} \\
Percept/observable & $\hatO_{\mathrm{cog}} = P_{\mathrm{cog}} \hatO P_{\mathrm{cog}}$ & Def.~\ref{def:16.1} \\
Cognitive dynamics & $\CL_{\mathrm{cog}}$ & Lemma~\ref{lem:16.1} \\
\bottomrule
\end{tabular}
\end{center}

\subsection{Mathematical Construction: Recursive Cognition}

\begin{definition}[Cognitive Recursion Operator]\label{def:16.4}
The \emph{cognitive recursion operator} $F_{\mathrm{cog}} :
\mathcal{O}_{\mathrm{cog}} \to \mathcal{O}_{\mathrm{cog}}$ is:
\[
  F_{\mathrm{cog}}(O_t) := P_{\mathrm{cog}}\, \Phi^*(O_t)\, P_{\mathrm{cog}}
  + \alpha\, P_{\mathrm{cog}}\, [\hatC, O_t]\, P_{\mathrm{cog}},
\]
where $\Phi^*$ is the adjoint of $\Phi$ (the pullback through the
interface map, Axiom~\ref{ax:4.5}), and $\alpha := \heff / \kappa_{\max}$.
\end{definition}

\begin{theorem}[Cognitive Recursion Dynamics]\label{thm:16.1}
Let $O_0 \in \mathcal{O}_{\mathrm{cog}}$ with $\|O_0\|_{\mathrm{op}} = 1$.
Define $O_{t+1} := F_{\mathrm{cog}}(O_t)$.  Then:
\begin{enumerate}
  \item[\textup{(i)}] The sequence $(O_t)_{t \geq 0}$ is bounded:
    $\|O_t\|_{\mathrm{op}} \leq (1 + \alpha \|\hatC\|_{\mathrm{op}})^t$.
  \item[\textup{(ii)}] If $\|\Phi^*\|_{\mathrm{op}} + \alpha \|\hatC\|_{\mathrm{op}} < 1$,
    the sequence converges to a fixed point $O^* \in \mathcal{O}_{\mathrm{cog}}$
    satisfying $F_{\mathrm{cog}}(O^*) = O^*$.
  \item[\textup{(iii)}] The fixed point $O^*$ represents the
    \emph{stable cognitive representation} of the external information
    mediated by $\Phi$.
\end{enumerate}
\end{theorem}

\begin{proof}
(i) follows from the submultiplicativity of the operator norm and
induction.  (ii) follows from the Banach fixed-point theorem applied
to $F_{\mathrm{cog}}$ on the complete metric space
$(\mathcal{O}_{\mathrm{cog}}, \|\cdot\|_{\mathrm{op}})$ with
contraction constant $c := \|\Phi^*\|_{\mathrm{op}} + \alpha
\|\hatC\|_{\mathrm{op}} < 1$.  (iii) is the interpretation:
$O^* = F_{\mathrm{cog}}(O^*)$ means the operator is invariant under
one step of interface-mediated perception plus constraint-induced
correction, i.e.\ a stable percept.
\end{proof}

\subsection{Derived Cognitive Constructs}

\begin{definition}[Attention Operator]\label{def:16.5}
The \emph{attention operator} $A_\omega$ for a spectral window
$\omega \subset \mathrm{spec}(\hatC)$ is:
\[
  A_\omega := \sum_{\lambda_j \in \omega} \hatP_{\lambda_j},
\]
where $\hatP_{\lambda_j}$ are the spectral projections of $\hatC$.
Attention is the act of restricting the cognitive state to a
subspectrum:
\[
  \rho \mapsto \rho_\omega :=
  \frac{A_\omega \rho A_\omega}{\Tr(A_\omega \rho A_\omega)}.
\]
\end{definition}

\begin{definition}[Memory Operator]\label{def:16.6}
The \emph{memory operator} $M_\tau$ for retention time $\tau > 0$
is the CPTP map:
\[
  M_\tau(\rho) := e^{\tau \CL_{\mathrm{cog}}}(\rho).
\]
Memory is the time-evolution of the cognitive state under the
domain-restricted Lindblad generator.  The Lindblad decay rates
$\gamma_k$ (Chapter~7) determine the forgetting timescales:
$\tau_k^{\mathrm{forget}} = \gamma_k^{-1}$.
\end{definition}

\begin{definition}[Perception Map]\label{def:16.7}
The \emph{perception map} $\Pi : \mathfrak{D}(\HC) \to
\Omega_{\mathrm{cog}}$ is:
\[
  \Pi(\rho) := \frac{P_{\mathrm{cog}} \,\Phi(\rho)\, P_{\mathrm{cog}}}
               {\Tr(P_{\mathrm{cog}} \,\Phi(\rho)\, P_{\mathrm{cog}})}.
\]
This is the composition of the interface map and the domain projection,
yielding the perceived state from an arbitrary density operator.
\end{definition}

\subsection{Derivation: Weber--Fechner Law}

\begin{proposition}[Constraint-Induced Weber--Fechner]\label{prop:16.1}
Let $\hatO_s$ be a stimulus-intensity observable with eigenvalue $s > 0$.
Under the spectral distortion bound (Theorem~8.2),
\[
  \langle \hatO_s \rangle_{\mathrm{cog}} =
  \langle P_{\mathrm{cog}} \hatO_s P_{\mathrm{cog}} \rangle_\rho
  = s + D(s),
\]
where $|D(s)| \leq C \kappa_{\max} s$ for a constant $C > 0$.
Then the \emph{perceived intensity} $\psi(s) :=
\langle \hatO_s \rangle_{\mathrm{cog}}$ satisfies, to first order
in $\kappa_{\max}$,
\[
  \frac{d\psi}{ds} = 1 + C \kappa_{\max} + O(\kappa_{\max}^2)
  \approx \frac{\mathrm{const}}{s}
\]
when the distortion $D(s) \sim \kappa_{\max} \ln s$ arises from
logarithmic dependence of the constraint curvature on stimulus
magnitude.  This yields:
\[
  \psi(s) \approx k \ln s + c_0
  \qquad (\text{Weber--Fechner law}).
\]
\end{proposition}

\subsection{Cross-Domain Links}

\begin{itemize}
  \item \textbf{To Physics (Ch.~14):} The cognitive constraint curvature
        $\kappa$ that determines the distortion bound is the same scalar
        curvature from which spacetime geometry emerges
        (Theorem~14.5).  Cognitive distortion is a finite-bandwidth
        restriction of the same information-geometric structure.
  \item \textbf{To Measurement Theory (Ch.~8):} The perception map
        $\Pi$ is the composition of the CIIR POVM
        (Definition~8.3) with the domain projection, inheriting all
        spectral distortion bounds.
\end{itemize}

% ============================================================
\section{Domain II: Neuroscience}
\label{sec:ch16:neuro}
% ============================================================

\subsection{Formal Definition}

\begin{definition}[Neural Domain]\label{def:16.8}
The \emph{neural domain} is:
\[
  D_{\mathrm{neur}} := (P_{\mathrm{neur}},\;
                        \mathcal{O}_{\mathrm{neur}},\;
                        \Omega_{\mathrm{neur}})
\]
where $P_{\mathrm{neur}}$ projects onto a tensor-product subspace
modelling $N$ interacting neural units:
\[
  \HC_{\mathrm{neur}} := \bigotimes_{i=1}^{N} \HC_i, \qquad
  \dim(\HC_i) = d_i < \infty.
\]
Each factor $\HC_i$ represents the state space of the $i$-th neural
unit; $d_i$ encodes the number of distinguishable activity levels.
\end{definition}

\subsection{Ontological Mapping}

\begin{center}
\begin{tabular}{lll}
\toprule
\textbf{Neural Concept} & \textbf{CIIR Object} & \textbf{Reference} \\
\midrule
Neural state & $\rho_{\mathrm{neur}} \in \Omega_{\mathrm{neur}}$ & Def.~\ref{def:16.8} \\
Synaptic weight & $W_{ij} \in \mathcal{O}_{\mathrm{neur}}$ & Def.~\ref{def:16.9} \\
Firing observable & $\sigma_i^z \otimes I_{\neq i}$ & --- \\
Neural dynamics & $\CL_{\mathrm{neur}}$ & Lemma~\ref{lem:16.1} \\
\bottomrule
\end{tabular}
\end{center}

\subsection{Mathematical Construction: Operator Networks}

\begin{definition}[Synaptic Operator]\label{def:16.9}
The \emph{synaptic operator} coupling neural units $i$ and $j$ is:
\[
  W_{ij} := w_{ij}\, \hatC_i \otimes \hatC_j \in \mathcal{O}_{\mathrm{neur}},
\]
where $w_{ij} \in \mathbb{R}$ is the coupling strength and
$\hatC_i$ is the constraint operator restricted to $\HC_i$.
The \emph{neural Hamiltonian} is:
\[
  \hatH_{\mathrm{neur}} := \sum_{i} h_i \hatC_i \otimes I_{\neq i}
  + \sum_{i < j} W_{ij}.
\]
\end{definition}

\begin{theorem}[Neural Network as Operator Network]\label{thm:16.2}
The neural domain $D_{\mathrm{neur}}$ with Hamiltonian
$\hatH_{\mathrm{neur}}$ is isomorphic (as a dynamical system) to
an operator network $\mathcal{N} = (\{O_i\}, \{W_{ij}\})$ where:
\begin{enumerate}
  \item[\textup{(i)}] Each node $O_i := \hatC_i$ is a constraint
    operator.
  \item[\textup{(ii)}] Each edge $W_{ij}$ is a tensor-product coupling
    operator.
  \item[\textup{(iii)}] The dynamics are generated by the Lindblad
    equation with Hamiltonian $\hatH_{\mathrm{neur}}$ and Lindblad
    operators $L_k^{(i)} := \sqrt{\gamma_k}\, \hatC_k^{(i)}$ for
    each unit $i$.
\end{enumerate}
The isomorphism is:
$\mathcal{N} \cong D_{\mathrm{neur}}$ as $C^*$-dynamical systems.
\end{theorem}

\begin{proof}
The mapping $O_i \mapsto \hatC_i \otimes I_{\neq i}$ and
$W_{ij} \mapsto w_{ij} \hatC_i \otimes \hatC_j$ is a faithful
$*$-homomorphism from the free operator network algebra to
$\mathcal{O}_{\mathrm{neur}}$.  The dynamics are preserved because
the Lindblad generator is constructed from the same operators.
\end{proof}

\subsection{Derivation: Hebb's Rule}

\begin{proposition}[Constraint-Induced Hebbian Learning]\label{prop:16.2}
Under the CIIR master equation restricted to $D_{\mathrm{neur}}$,
the expectation of the synaptic operator evolves as:
\[
  \frac{d}{dt}\langle W_{ij} \rangle
  = w_{ij} \Bigl(
    \langle \hatC_i \rangle \langle \hatC_j \rangle
    + \mathrm{Cov}(\hatC_i, \hatC_j)
    - \gamma_{\mathrm{decay}} \langle W_{ij} \rangle
  \Bigr),
\]
where $\mathrm{Cov}(\hatC_i, \hatC_j) :=
\langle \hatC_i \hatC_j \rangle - \langle \hatC_i \rangle
\langle \hatC_j \rangle$.  In the mean-field limit (negligible
covariance), this reduces to:
\[
  \frac{d}{dt} w_{ij}^{\mathrm{eff}} \approx
  \eta \langle \hatC_i \rangle \langle \hatC_j \rangle
  - \gamma_{\mathrm{decay}}\, w_{ij}^{\mathrm{eff}}
  \qquad (\text{Hebb's rule with decay}),
\]
where $\eta$ is an effective learning rate determined by the constraint
curvature.
\end{proposition}

\subsection{Cross-Domain Links}

\begin{itemize}
  \item \textbf{To Cognitive Science (Sec.~\ref{sec:ch16:cogsci}):}
        The cognitive domain projection $P_{\mathrm{cog}}$ is the
        trace over internal neural degrees of freedom:
        $\rho_{\mathrm{cog}} = \Tr_{\mathrm{int}}(\rho_{\mathrm{neur}})$.
  \item \textbf{To Physics (Ch.~14):} The tensor-product structure of
        $\HC_{\mathrm{neur}}$ mirrors the multi-particle Hilbert space;
        entanglement between neural units is constraint-induced
        entanglement (Theorem~10.3).
\end{itemize}

% ============================================================
\section{Domain III: Interface Theory and Human--Computer Interaction}
\label{sec:ch16:hci}
% ============================================================

\subsection{Formal Definition}

\begin{definition}[Interface Domain]\label{def:16.10}
The \emph{interface domain} is:
\[
  D_{\mathrm{HCI}} := (P_{\mathrm{HCI}},\;
                       \mathcal{O}_{\mathrm{HCI}},\;
                       \Omega_{\mathrm{HCI}})
\]
where $P_{\mathrm{HCI}}$ projects onto the subspace accessible through
the interface map $\Phi$:
\[
  \HC_{\mathrm{HCI}} := \overline{\mathrm{Im}(\Phi)}
  = \HC_\CM \quad (\text{Definition~\ref{def:3.21}}).
\]
\end{definition}

\begin{remark}
The HCI domain is the constraint-image subspace $\HC_\CM$ itself.
This means that the interface map $\Phi$ is, by construction, the
formalisation of the human--computer (or agent--reality) interface.
All of CIIR's spectral distortion, non-invertibility, and contextuality
results (Chapter~8) are natively HCI results.
\end{remark}

\subsection{Theorem: Fitness-Based Interface}

\begin{theorem}[Interface Optimality Under Constraints]\label{thm:16.3}
Let $\Phi$ be the interface map and $\hatC$ the constraint operator.
Among all linear maps $\Psi : \CB(\CR_{\mathrm{op}}) \to \CB(\HC)$
satisfying Axioms~\ref{ax:4.4}--\ref{ax:4.5}, the interface map
$\Phi$ minimises the constraint functional:
\[
  \CC_{\mathrm{HCI}}(\Psi) := \Tr\bigl(\hatC\, \Psi(\hatC)\bigr)
  + \mu \|\Psi - \Phi\|_{\mathrm{op}}^2,
\]
for a regularisation parameter $\mu > 0$ determined by
$\kappa_{\max}$.  That is, $\Phi$ is the \emph{fitness-optimal}
interface in the sense of Hoffman's interface theory of perception.
\end{theorem}

\subsection{Cross-Domain Links}

\begin{itemize}
  \item \textbf{To Cognitive Science (Sec.~\ref{sec:ch16:cogsci}):}
        The perception map $\Pi$ (Definition~\ref{def:16.7}) is the
        composition $\Pi = P_{\mathrm{cog}} \circ \Phi$, linking HCI
        directly to cognition.
  \item \textbf{To Measurement (Ch.~8):} All POVM results apply
        directly to the HCI domain since $\HC_{\mathrm{HCI}} = \HC_\CM$.
\end{itemize}

% ============================================================
\section{Domain IV: Systems Biology and Evolutionary Biology}
\label{sec:ch16:biology}
% ============================================================

\subsection{Formal Definition}

\begin{definition}[Biological Domain]\label{def:16.11}
The \emph{biological domain} is:
\[
  D_{\mathrm{bio}} := (P_{\mathrm{bio}},\;
                       \mathcal{O}_{\mathrm{bio}},\;
                       \Omega_{\mathrm{bio}})
\]
where $P_{\mathrm{bio}}$ projects onto a subspace
$\HC_{\mathrm{bio}} \subseteq \HC_{\mathrm{neur}}$ encoding
organism-level configurations.  The \emph{fitness functional}
$\mathcal{F} : \Omega_{\mathrm{bio}} \to \mathbb{R}_{\geq 0}$
is:
\[
  \mathcal{F}(\rho_{\mathrm{bio}}) :=
  -\Tr\bigl(\hatC_{\mathrm{bio}}\, \rho_{\mathrm{bio}}\bigr),
\]
where $\hatC_{\mathrm{bio}} := P_{\mathrm{bio}} \hatC P_{\mathrm{bio}}$.
Higher fitness corresponds to lower expected constraint.
\end{definition}

\subsection{Mathematical Construction: Evolution as Constraint Gradient Flow}

\begin{definition}[Evolutionary Flow]\label{def:16.12}
The \emph{evolutionary flow} on $\Omega_{\mathrm{bio}}$ is the
gradient flow of the fitness functional with respect to the
Fisher--Rao metric (Theorem~14.1):
\[
  \frac{\partial \rho_{\mathrm{bio}}}{\partial t}
  = -\nabla_{g_F} \CC(\rho_{\mathrm{bio}}),
\]
where $g_F$ is the Fisher information metric on $\Omega_{\mathrm{bio}}$.
\end{definition}

\begin{theorem}[Natural Selection as Constraint Minimisation]\label{thm:16.4}
The evolutionary flow of Definition~\ref{def:16.12} converges to a
state $\rho^*_{\mathrm{bio}}$ that minimises the constraint functional
$\CC$ restricted to $\Omega_{\mathrm{bio}}$, subject to the topological
constraints of the fitness landscape.  In the finite-dimensional case
$(\dim \HC_{\mathrm{bio}} = d < \infty)$:
\begin{enumerate}
  \item[\textup{(i)}] The flow exists globally and is unique.
  \item[\textup{(ii)}] Every trajectory converges to a critical point
    of $\CC|_{\Omega_{\mathrm{bio}}}$.
  \item[\textup{(iii)}] If $\CC$ is strictly convex on
    $\Omega_{\mathrm{bio}}$, the minimiser $\rho^*_{\mathrm{bio}}$ is
    unique.
\end{enumerate}
This identifies natural selection with constraint minimisation on
the informational manifold.
\end{theorem}

\begin{proof}
The Fisher--Rao metric is positive definite on
$\Omega_{\mathrm{bio}}$ (a finite-dimensional statistical manifold).
The constraint functional $\CC$ is continuous and bounded below on
the compact set $\Omega_{\mathrm{bio}}$ (since density operators form
a compact convex set in finite dimensions).  Standard gradient-flow
theory on Riemannian manifolds (cf.\ Ambrosio et al.) gives global
existence, uniqueness, and convergence to critical points.  Strict
convexity gives uniqueness of the minimiser.
\end{proof}

\subsection{Cross-Domain Links}

\begin{itemize}
  \item \textbf{To Neuroscience (Sec.~\ref{sec:ch16:neuro}):}
        The biological domain is a coarse-graining of the neural domain:
        $\HC_{\mathrm{bio}} \subset \HC_{\mathrm{neur}}$, and the fitness
        functional is the partial trace of the constraint operator.
  \item \textbf{To Emergent Physics (Ch.~14):} The Fisher--Rao metric
        used in the evolutionary flow is the same information metric
        from which spacetime geometry emerges (Theorem~14.1),
        establishing a direct structural link between evolution and
        geometry.
\end{itemize}
